% ЛР №3. Содержание перенесено из пользовательского DOCX без редакции.
% Компилятор: XeLaTeX. Единый согласованный шаблон mietlab.
\documentclass{mietlab}
\LabSetup{
  number=3,
  title={«Интеграционное тестирование»},
  variant={Вариант 3},
  student={Варлапов Артем},
  group={П-42},
  studentline={студент группы П-42},
  institute={Институт системной и программной инженерии и информационных технологий},
  discipline={Основы программной инженерии},
  city={Москва},
  year={2026}
}
\begin{document}
\MakeLabTitle
\LabHeading{Цель работы}

Научиться выполнять интеграционное тестирование ПО.\par

\LabHeading{Задание}

Сформировать иерархию классов вокруг ранее протестированного класса. Провести модульное тестирование написанных классов и инкрементальное интеграционное тестирование сверху вниз с использованием программных заглушек для имитации ещё не подключённых классов нижнего уровня.\par

Для варианта 3 исходным является класс «Строка» TextLine из лабораторной работы №2. В лабораторной работе №3 отдельного списка вариантов нет, поэтому ссылка на класс из лабораторной №3 трактуется как ссылка на ранее протестированный класс из лабораторной №2.\par

\LabHeading{Выполнение работы}

\LabHeading{1 Иерархия и взаимодействие классов}

\begin{sloppypar}
TextLine является базовым классом строки. Классы RussianTextLine и EnglishTextLine наследуют TextLine и задают язык при создании объекта. TextDocument является агрегирующим классом: он хранит список объектов TextLine и суммирует их текущие длины.\par
\end{sloppypar}

Связи классов: RussianTextLine → TextLine и EnglishTextLine → TextLine являются наследованием; TextDocument → список TextLine является агрегацией. При вычислении длины документа вызывается getLength() каждого объекта строки.\par

\begin{labtable}{|L{140bp}|L{\dimexpr\linewidth-140bp-\arrayrulewidth\relax}|}
\hline
\LabTableHead{Класс} & \LabTableHead{Назначение} \\\hline
\endfirsthead
\hline
\LabTableHead{Класс} & \LabTableHead{Назначение} \\\hline
\endhead
TextLine & Базовый класс строки: хранит массив символов, длину и язык; поддерживает изменение длины и конкатенацию. \\\hline
RussianTextLine & Наследник TextLine, задающий язык «Русский» при создании объекта. \\\hline
EnglishTextLine & Наследник TextLine, задающий язык «Английский» при создании объекта. \\\hline
TextDocument & Агрегирующий класс, хранящий список строк и вычисляющий суммарную длину документа. \\\hline
TextLineStub & Программная заглушка для верхнеуровневого тестирования TextDocument. \\\hline
\end{labtable}

\LabHeading{2 Реализация классов}

Базовый класс TextLine перенесён из лабораторной работы №2. Некорректные аргументы обрабатываются через IllegalArgumentException. Наследники задают фиксированный язык, а TextDocument добавляет проверку null при добавлении строки.\par

\Needspace{156bp}
\LabHeading{TextLine.java}

\LabCodeFile[language=Java]{src/TextLine.java}

\Needspace{156bp}
\LabHeading{RussianTextLine.java}

\LabCodeFile[language=Java]{src/RussianTextLine.java}

\Needspace{156bp}
\LabHeading{EnglishTextLine.java}

\LabCodeFile[language=Java]{src/EnglishTextLine.java}

\Needspace{156bp}
\LabHeading{TextDocument.java}

\LabCodeFile[language=Java]{src/TextDocument.java}

\LabHeading{3 Заглушка и тестирование сверху вниз}

На первом этапе проверяется верхний класс TextDocument. Настоящие строки заменяются объектами TextLineStub. Переопределённый метод getLength() возвращает заданное тестом число независимо от внутреннего массива символов. Наследование используется для совместимости с интерфейсом TextLine.\par

\Needspace{156bp}
\LabHeading{TextLineStub.java}

\LabCodeFile[language=Java]{src/TextLineStub.java}

Этап 1: TextDocumentTest проверяет пустой документ, одну заглушку, сумму длин нескольких заглушек и отклонение null. Так проверяется логика документа до подключения настоящих строк.\par

Этап 2: в тесте documentWorksWithRealTextLines заглушки заменяются настоящими объектами TextLine. Проверяется сумма длин строк «abc» и «de», равная 5 символам.\par

\begin{sloppypar}
Этап 3: в тестах documentWorksWithBothSubclasses и documentReflectsChangesInItsLines подключаются наследники. Проверяется сумма длин русской и английской строк, а также изменение длины документа после append() и setLength() у сохранённой строки.\par
\end{sloppypar}

\Needspace{475bp}
\LabHeading{4 Модульные тесты}

Написаны тесты всех четырёх классов программы. Всего подготовлено 24 модульных теста. Проверяются корректные данные, пустые строки, граничные длины и ошибочные аргументы.\par

\begin{labtable}{|L{175.5bp}|L{95bp}|L{\dimexpr\linewidth-270.5bp-\arrayrulewidth\relax}|}
\hline
\LabTableHead{Тестовый класс} & \LabTableHead{Количество тестов} & \LabTableHead{Что проверяется} \\\hline
\endfirsthead
\hline
\LabTableHead{Тестовый класс} & \LabTableHead{Количество тестов} & \LabTableHead{Что проверяется} \\\hline
\endhead
TextLineTest & 14 & Проверка конструктора, изменения языка, длины, конкатенации и ошибочных аргументов. \\\hline
RussianTextLineTest & 3 & Проверка языка, длины пустой строки и обработки null. \\\hline
EnglishTextLineTest & 3 & Проверка языка, длины пустой строки и обработки null. \\\hline
TextDocumentTest & 4 & Верхнеуровневая проверка документа с программными заглушками. \\\hline
TextDocumentIntegrationTest & 3 & Проверка взаимодействия документа с настоящими строками и наследниками. \\\hline
\end{labtable}

\Needspace{156bp}
\LabHeading{TextLineTest.java}

\LabCodeFile[language=Java]{src/TextLineTest.java}

\Needspace{156bp}
\LabHeading{RussianTextLineTest.java}

\LabCodeFile[language=Java]{src/RussianTextLineTest.java}

\Needspace{156bp}
\LabHeading{EnglishTextLineTest.java}

\LabCodeFile[language=Java]{src/EnglishTextLineTest.java}

\Needspace{156bp}
\LabHeading{TextDocumentTest.java}

\LabCodeFile[language=Java]{src/TextDocumentTest.java}

\LabHeading{5 Интеграционные тесты}

Три интеграционных теста проверяют взаимодействие документа с настоящим базовым классом и наследниками, включая пересчёт общей длины после изменения строки.\par

\Needspace{156bp}
\LabHeading{TextDocumentIntegrationTest.java}

\LabCodeFile[language=Java]{src/TextDocumentIntegrationTest.java}

\Needspace{520bp}
\LabHeading{6 Результаты тестирования}

Сначала отдельно запущена проверка верхнего класса с заглушками. Все 4 теста прошли успешно.\par

\Needspace{117bp}
\LabHeading{Результат проверки TextDocument с программными заглушками}

\LabCodeFile{results/stub-tests.txt}

Затем запущен полный набор модульных и интеграционных тестов через JUnitCore из JUnit 4.13.2. Все 27 тестов прошли успешно.\par

\Needspace{117bp}
\LabHeading{Результат выполнения полного набора тестов}

\LabCodeFile{results/all-tests.txt}

\LabHeading{Вывод}

Удалось научиться выполнять интеграционное тестирование ПО\par

\end{document}
